\begin{env}[3.8.3]
\label{0.3.8.3-fr}
Sean $\mathcal{F}$ un haz de conjuntos, $\mathcal{H}$ un subhaz de
$\mathcal{F}$, y $\mathcal{F}'$ y $\mathcal{H}'$ los haces pseudodiscretos
asociados a $\mathcal{F}$ y $\mathcal{H}$, respectivamente. Entonces, para
todo abierto $U\subset X$, $\Gamma(U,\mathcal{H}')$ es \emph{cerrado} en
$\Gamma(U,\mathcal{F}')$. En efecto, es la intersección de las imágenes
inversas de los $\Gamma(V,\mathcal{H})$ (para $V\in\mathfrak{B}$,
$V\subset U$) por las aplicaciones continuas
$\Gamma(U,\mathcal{F})\to\Gamma(V,\mathcal{F})$, y
$\Gamma(V,\mathcal{H})$ es cerrado en el espacio discreto
$\Gamma(V,\mathcal{F})$.
\end{env}
